% The sharp constant in the shift-wall bridge for odd-cycle polynomials.
% Author: J. Councilman
% Compile with: latexmk -pdf shift_wall_sharp_factor.tex

\documentclass[11pt,a4paper]{article}

\usepackage[margin=1.1in]{geometry}
\usepackage{amsmath,amssymb,amsthm}
\usepackage{array,booktabs,tabularx}
\usepackage{enumitem}
\usepackage[colorlinks=true,linkcolor=black,citecolor=black,urlcolor=blue]{hyperref}
\hypersetup{
  pdftitle={The sharp constant in the shift-wall bridge for odd-cycle polynomials},
  pdfauthor={J. Councilman},
  pdfsubject={The exact extremal constant in a Turan bridge between consecutive wall rows},
  pdfkeywords={odd cycles; Wall recurrence; Turan inequality; shift wall;
               exact certificates; sharp constant}
}
\usepackage[expansion=false]{microtype}
\setlength{\emergencystretch}{2em}
\widowpenalty=10000
\clubpenalty=10000
\displaywidowpenalty=10000

\theoremstyle{plain}
\newtheorem{theorem}{Theorem}[section]
\newtheorem{proposition}[theorem]{Proposition}
\newtheorem{lemma}[theorem]{Lemma}
\newtheorem{corollary}[theorem]{Corollary}
\newtheorem{conjecture}[theorem]{Conjecture}
\theoremstyle{remark}
\newtheorem{remark}[theorem]{Remark}

\newcommand{\N}{\mathbb{N}}
\newcommand{\Q}{\mathbb{Q}}
\newcommand{\Omw}{\omega_{\mathrm{wall}}}
\newcommand{\Kw}{K_{\mathrm{wall}}}
\newcommand{\level}{\mathfrak a}
\newcolumntype{P}[1]{>{\raggedright\arraybackslash}p{#1}}
\newcolumntype{Y}{>{\raggedright\arraybackslash}X}

\title{\textbf{The sharp constant in the shift-wall bridge\\
for odd-cycle polynomials}}
\author{J. Councilman\\
\small\href{https://orcid.org/0009-0003-6265-5168}{ORCID 0009-0003-6265-5168}}
\date{9 August 2026}

\begin{document}
\maketitle

\begin{abstract}
Let
\[
 M_0=1,\qquad M_1=1+z,\qquad
 M_n=M_{n-1}+n^2zM_{n-2},
\]
and write $L(f)_q=f_q^2-f_{q-1}f_{q+1}$ for the coefficientwise
Tur\'an operator.  A previous square-tail theorem implies
\[
 (2r)^2L\bigl((1+z)M_{2r-2}\bigr)_q\ge 2L(M_{2r-1})_q,
 \qquad 0\le q\le r-2.
\]
We determine the best possible replacement for the factor $2$.  Put
\[
 \Omega_t(r)=
 \frac{(2r)^2L((1+z)M_{2r-2})_{r-t}}
      {2L(M_{2r-1})_{r-t}},\qquad r\ge t\ge2.
\]
Then
\[
 \Omega_t(r)\ge \Omega_2(1350)
 =1.327509156956680599863203087838\ldots,
\]
with equality only at $(r,t)=(1350,2)$.  Consequently the sharp constant in
the unnormalized bridge is
\[
 2\Omega_2(1350)
 =2.655018313913361199726406175676\ldots.
\]
The proof has one exceptional stripe and one uniform barrier.  The $t=2$
minimum is certified by an exact prefix and a positive-coefficient half-line
certificate.  A lossless telescoping transfer sends every fixed-column lower
bound for the square-tail quotient $\rho$ to the corresponding bridge stripe.
The stripe $t=3$ is treated directly, while $\rho>4/3$ closes every stripe
$t\ge4$.  For $t\ge5$ this barrier is assembled from terminal identities, a
$123{,}753$-cell exact core, six fixed-column tail certificates, a re-anchored
top estimate, and an all-parameter joint certificate; on that region it is
also a uniform improvement, by a factor $1.314\ldots$, on the sharp
square-tail constant of the earlier paper.  A byproduct localizes the whole
question: every cell below $4/3$ lies on the stripe $t=2$ with $r\le3000$.
Every decision gate is an integer or rational comparison; decimal values are
reported only for orientation.  The stripes $t=9$ and $t=10$ each carry a
second, separately gated proof chain, one of which rests on an exact two-leg
certificate for a quantitative decrement shell.
\end{abstract}

\section{The extremal question}\label{sec:intro}

For a polynomial $f(z)=\sum_k f_kz^k$, put
\begin{equation}\label{eq:L}
 L(f)_q=f_q^2-f_{q-1}f_{q+1},
\end{equation}
where coefficients outside the row are zero.  The wall recurrence
\begin{equation}\label{eq:wall}
 M_0=1,\qquad M_1=1+z,\qquad
 M_n=M_{n-1}+n^2zM_{n-2}
\end{equation}
is the even-variable form of the odd-cycle continuant.  The rows are
real-rooted with positive coefficients, so the Tur\'an forms used below are
positive on their natural range; the determinant, spectral, and permutation
interpretations are developed in~\cite{squaretail}.

For $r\ge2$ and $0\le q\le r-2$, define the raw bridge ratio
\[
 \mathcal B(r,q)=
 \frac{(2r)^2L((1+z)M_{2r-2})_q}{L(M_{2r-1})_q}.
\]
Writing $t=r-q$ reverses the coefficient index and gives
\begin{equation}\label{eq:Omega}
 \Omega_t(r)=\frac{\mathcal B(r,r-t)}2
 =\frac{(2r)^2L((1+z)M_{2r-2})_{r-t}}
        {2L(M_{2r-1})_{r-t}},
 \qquad r\ge t\ge2.
\end{equation}
The bridge theorem in~\cite{squaretail} is $\Omega_t(r)\ge1$.  Its constant
$2$ is structural: it is the amount paid by the square-tail inequality and
the Tur\'an triangle.  The present question is extremal: what is
$\inf_{r\ge t\ge2}\Omega_t(r)$?

Set
\begin{equation}\label{eq:constants}
 \Omw=\Omega_2(1350),\qquad \Kw=2\Omw.
\end{equation}

\begin{theorem}[sharp shift-wall constant]\label{thm:main}
For all integers $r\ge t\ge2$,
\begin{equation}\label{eq:main}
 \Omega_t(r)\ge\Omw
 =1.327509156956680599863203087838\ldots,
\end{equation}
with equality if and only if $(r,t)=(1350,2)$.  Equivalently,
\begin{equation}\label{eq:rawsharp}
 (2r)^2L\bigl((1+z)M_{2r-2}\bigr)_q
 \ge \Kw L(M_{2r-1})_q,
 \qquad 0\le q\le r-2,
\end{equation}
where
\[
 \Kw=2.655018313913361199726406175676\ldots.
\]
The constant is best possible, and equality in \eqref{eq:rawsharp} occurs
only at $(r,q)=(1350,1348)$.
\end{theorem}

The exact value in \eqref{eq:main} is a reduced rational.  Its numerator has
$6{,}733$ decimal digits, its denominator $6{,}732$, and the SHA-256 digest of
the canonical ASCII string \texttt{numerator/denominator} is
\begin{center}
\texttt{2fc6689fd42dd6eaa29088aeb0077e1120878fd2f6b54cf2e8dbd50e823a5691}.
\end{center}
The decimal is not used as a proof target.

The proof is organized around the rational separator
\begin{equation}\label{eq:separator}
 \Omw<\frac43,
 \qquad
 \frac43-\Omw
 =0.005824176376652733470130\ldots.
\end{equation}
The exceptional stripe $t=2$ supplies the sharp value.  Every stripe $t\ge3$
is strictly above $4/3$.  The main work is therefore not a decimal comparison
with \eqref{eq:main}, but the uniform exact theorem $\rho>4/3$ proved in
Section~\ref{sec:bulk}.

\subsection{Place in the four-paper program}

The logical input is the square-tail theorem of~\cite{squaretail}.  That paper
builds the odd-cycle, Jacobi, and coefficient-curvature machinery and proves
the bridge with the structural factor $2$; this paper determines the exact
constant left open there.  It also left a prediction: its own sharp
square-tail constant $C_\star=2\rho(129,2)=2.0291523\ldots$ propagates
through the bridge, but is not sharp there, and the bridge's extremal cell is
\emph{not} the image of $(129,2)$.  Theorem~\ref{thm:main} confirms both.  The
extremal cell is $(1350,2)$, and $\Kw$ exceeds $C_\star$ by $30.8\%$ and the
structural factor $2$ by $32.8\%$.  The sharp multiplicative Delannoy theorem
in~\cite{delannoy} is the external application that first produced the bridge
question.  The grounded-path algorithms of~\cite{greenpath} expose a related
determinantal mechanism, but are a conceptual sibling rather than a premise of
Theorem~\ref{thm:main}.

\section{The square-tail quotient and the wall level}\label{sec:objects}

Reverse and normalize the even wall row by
\begin{equation}\label{eq:Pi}
 \Pi_r(u)=\frac{u^{r-1}M_{2r-2}(1/u)}{(2r-1)!}.
\end{equation}
For $r\ge2$ and $2\le t\le r$, define
\begin{align}
 S_{r,t}&=L\bigl((1+u)\Pi_r\bigr)_t,&
 V_{r,t}&=L(\Pi_r)_{t-1},\label{eq:SV}\\
 d_{r,t}&=S_{r,t}-\left(\frac{2r-1}{2r}\right)^2S_{r-1,t},&
 \rho(r,t)&=\frac{r^2d_{r,t}^2}{2S_{r,t}V_{r,t}}.\label{eq:drho}
\end{align}
The results of~\cite{squaretail} give $S_{r,t},V_{r,t},d_{r,t}>0$ and
$\rho(r,t)\ge1$ on this domain.  The quotient $\rho$ measures the exact
surplus in the cross-row square-tail inequality.  It is not the bridge ratio
$\Omega$: the next section records the one-way transfer between them.

The unbounded part of the $\rho$-array is organized by one increasing scalar.
Let
\[
 \widehat\Pi_r(u)=\frac{\Pi_r(u)}{\Pi_r(0)},
 \qquad
 \level_r=6[u]\widehat\Pi_r(u).
\]
The cancellation-free Wallis ladder in~\cite{squaretail} proves that
$\level_r$ is strictly increasing.  Its role is geometric.  For fixed $t$,
the inequality
\begin{equation}\label{eq:topcondition}
 \level_r\le(t-2)^2
\end{equation}
defines the top region; its complement is the joint region.  The proof never
uses a floating approximation to decide which side of
\eqref{eq:topcondition} contains a cell.  The few crossings needed below are
certified by outward rational intervals.

\section{A lossless fixed-column transfer}\label{sec:transfer}

The bridge in~\cite{squaretail} was obtained by inserting the square-tail
constant $1$ at every row.  Keeping the local quotient $\rho(j,t)$ instead
produces a sharper inequality with no additional loss.

For $0\le t\le r$, put
\[
 O_{r,t}=[z^{r-t}]M_{2r-1},\qquad
 E_{r,t}=[z^{r-t}]M_{2r},
\]
and normalize
\[
 \mathsf P_r(u)=\sum_t\frac{O_{r,t}}{O_{r,0}}u^t,
 \qquad
 \mathsf Q_r(u)=\sum_t\frac{E_{r,t}}{E_{r,0}}u^t,
 \qquad
 x_r=\frac{E_{r-1,0}}{O_{r,0}}.
\]
The wall recurrence gives
\begin{equation}\label{eq:PQrec}
 \mathsf P_r=\mathsf P_{r-1}+ux_r\mathsf Q_{r-1}.
\end{equation}
For a fixed $t\ge2$, set
\begin{align*}
 p_j&=\sqrt{L(\mathsf P_j)_t},\\
 A_j&=jx_j\sqrt{L((1+u)\mathsf Q_{j-1})_t},\\
 v_j&=x_j\sqrt{L(\mathsf Q_{j-1})_{t-1}}.
\end{align*}
All three depend on $t$, which is suppressed from the notation.

\begin{lemma}[lossless transfer]\label{lem:transfer}
For every pair of integers $r\ge t\ge2$, there are positive weights
\[
 \lambda_{j,t}=\frac{A_j-A_{j-1}}{A_r},
 \qquad t\le j\le r,
 \qquad \sum_{j=t}^r\lambda_{j,t}=1,
\]
such that
\begin{equation}\label{eq:transfer}
 \Omega_t(r)\ge
 \left(\sum_{j=t}^r
 \frac{\lambda_{j,t}}{\sqrt{\rho(j,t)}}\right)^{-2}
 \ge \min_{t\le j\le r}\rho(j,t).
\end{equation}
\end{lemma}

\begin{proof}
Two facts from~\cite{squaretail} carry the argument.  First, if $f$ and $g$
are real-rooted with nonnegative coefficients and lie in \emph{proper
position}, then polarizing the Newton inequality for the pencil $f+vg$ gives
the Tur\'an triangle
\begin{equation}\label{eq:turan}
 \sqrt{L(f+g)_t}\le\sqrt{L(f)_t}+\sqrt{L(g)_t}.
\end{equation}
The position hypothesis is not removable, and the minimal witness is
integral: $f=(1,2,1)$ and $g=(1,3,1)$ are both real-rooted with positive
coefficients, yet at $t=1$ they give $\sqrt{21}>\sqrt3+\sqrt8$, which on
squaring twice is the integer statement $100>96$.  What fails is position
alone --- the double root of $f$ at $-1$ lies strictly between the two roots
of $g$.  Second, the two rows appearing in
\eqref{eq:PQrec} do lie in proper position.  The zeros of $\mathsf P_{j-1}$
are $-\lambda_i^2$ and those of $u\mathsf Q_{j-1}$ are the $-\nu_i^2$
together with $0$, where $\lambda_i$ and $\nu_i$ run over the positive
spectra of the two consecutive Jacobi matrices of~\cite{squaretail}; strict
Cauchy interlacing and spectral symmetry give
\[
 0<\lambda_1<\nu_1<\cdots<\lambda_{j-1}<\nu_{j-1},
\]
so the zeros strictly interlace.  Applying \eqref{eq:turan} to
\eqref{eq:PQrec} therefore gives the Tur\'an triangle
\begin{equation}\label{eq:triangle}
 p_j\le p_{j-1}+v_j.
\end{equation}
In the same normalization, $A_j^2-A_{j-1}^2$ is the oriented decrement and
\[
 \rho(j,t)=\frac{(A_j^2-A_{j-1}^2)^2}{2A_j^2v_j^2}.
\]
The decrement is positive, so
\begin{equation}\label{eq:increment}
 A_j-A_{j-1}
 =\frac{A_j^2-A_{j-1}^2}{A_j+A_{j-1}}
 \ge \sqrt{\frac{\rho(j,t)}2}\,v_j.
\end{equation}
Both $p_{t-1}$ and $A_{t-1}$ vanish.  Iterating \eqref{eq:triangle}, using
\eqref{eq:increment}, and then telescoping gives
\[
 p_r\le\sum_{j=t}^rv_j
 \le\sqrt2\sum_{j=t}^r
 \frac{A_j-A_{j-1}}{\sqrt{\rho(j,t)}}
 =\sqrt2A_r\sum_{j=t}^r
 \frac{\lambda_{j,t}}{\sqrt{\rho(j,t)}}.
\]
Reversal in \eqref{eq:Omega} is exactly
$\Omega_t(r)=2A_r^2/p_r^2$.  This proves the first inequality in
\eqref{eq:transfer}.  The second follows because the $\lambda_{j,t}$ are
positive and sum to one.
\end{proof}

Thus a lower bound valid on an entire $\rho$-column transfers to the bridge
at the same constant.  The direction is important: no monotonicity of
$\Omega_t(r)$ in $t$ is asserted or needed.

\section{The exceptional stripe}\label{sec:t2}

The sharp value occurs on $t=2$, so this stripe is proved directly from the
bridge ratio rather than through Lemma~\ref{lem:transfer}.

Put $m=r-1$ and introduce the Wallis ladder of~\cite{squaretail},
\[
 h_j=\left(\frac{(2j)!!}{(2j-1)!!}\right)^2,
 \qquad
 T_j=1+\sum_{i=1}^{j}h_i^{-1},
 \qquad
 v_m=\frac{h_m}{(2m+1)^2},
\]
whose rows are determined by $B_0(j)=T_j$, $C_0(j)=1$,
\begin{equation}\label{eq:ladder}
 B_k(j)=B_k(j-1)+h_j^{-1}C_k(j-1),
 \qquad
 C_k(j)=C_k(j-1)+\frac{h_j}{(2j+1)^2}B_{k-1}(j),
\end{equation}
with seeds $B_k(0)=0$ for $k\ge1$, $C_1(0)=1$, and $C_k(0)=0$ for $k\ge2$.

The ladder is a rewriting of the two normalized rows of
\S\ref{sec:transfer}.  On the even side this is proved in~\cite{squaretail},
in the equivalent form $B_k(j)/T_j=[u^k]\Pi_{j+1}/[u^0]\Pi_{j+1}$:
\begin{equation}\label{eq:Bid}
 \frac{B_k(j)}{T_j}=[u^k]\mathsf Q_j .
\end{equation}
The odd side is the companion statement.  We record it because it is what
makes \eqref{eq:reverseOmega} a rewriting rather than a numerical
coincidence:
\begin{equation}\label{eq:Cid}
 C_k(m)=[u^k]\mathsf P_r
 =\frac{[z^{\,r-k}]M_{2r-1}}{[z^{\,r}]M_{2r-1}} .
\end{equation}
For \eqref{eq:Cid}, the wall recurrence first gives
\begin{equation}\label{eq:xladder}
 x_{j+1}=\frac{h_jT_j}{(2j+1)^2}.
\end{equation}
Now $C_0(j)=1=[u^0]\mathsf P_{j+1}$, and the coefficient of $u^k$ in
\eqref{eq:PQrec} reads
\[
 [u^k]\mathsf P_{j+1}=[u^k]\mathsf P_j+x_{j+1}[u^{k-1}]\mathsf Q_j .
\]
By \eqref{eq:Bid} and \eqref{eq:xladder} the last term equals
$h_j(2j+1)^{-2}B_{k-1}(j)$, so this recursion is exactly the second half of
\eqref{eq:ladder}.  Induction on $j$ proves \eqref{eq:Cid}.

Substituting \eqref{eq:Bid} and \eqref{eq:Cid} into \eqref{eq:Omega} and
cancelling the common normalizations gives
\begin{equation}\label{eq:reverseOmega}
 \Omega_t(r)=2r^2v_m^2\frac{P_t(B(m))}{Q_t(C(m))},
\end{equation}
where
\begin{align*}
 P_t(B)&=(B_{t-1}+B_t)^2
 -(B_{t-2}+B_{t-1})(B_t+B_{t+1}),\\
 Q_t(C)&=C_t^2-C_{t-1}C_{t+1}.
\end{align*}
Equation~\eqref{eq:reverseOmega} was also verified directly against the raw
wall recurrence in exact rational arithmetic at every cell with
$2\le t\le r\le9$, and against the rational ladder at the extremal row.

\begin{proposition}[the $t=2$ minimum]\label{prop:t2}
For every integer $r\ge2$,
\[
 \Omega_2(r)\ge\Omega_2(1350),
\]
with equality only at $r=1350$.
\end{proposition}

\begin{proof}[Exact certificate]
The prefix computes all $3{,}000$ cells $2\le r\le3001$ as reduced rational
comparisons with $\Omega_2(1350)$.  Every numerator and denominator is
positive, the equality set is $\{1350\}$, and all other cells are strictly
larger.  The nearest competitor is $r=1349$, whose positive gap is
\[
 3.473032420463471\ldots\times10^{-9}.
\]
The comparison stream has SHA-256 digest
\begin{center}
\texttt{943bebfb449b9ec178562598433b1f76dd9f3cb0cab0e8203bfce784882245ae}.
\end{center}

For the half-line, set $b_k=(\pi/2)B_k$ and
$x=\tau_m-\tau_{3000}^{\rm lo}\ge0$.  Outward rational envelopes give
\[
 P_2(b)\ge P_{\rm lo}(x),
 \qquad
 0<Q_{\rm lo}(x)\le Q_2(C)\le Q_{\rm hi}(x).
\]
The exact Wallis bound supplies a strict scalar factor $k_m>1/2$.  The
degree-ten polynomial
\[
 H(x)=P_{\rm lo}(x)-2\Omega_2(1350)Q_{\rm hi}(x)
\]
has all eleven rational coefficients strictly positive.  Its indexed
coefficient digest is
\begin{center}
\texttt{d8ba995d1f52b231ff5e198bd881ec9d1bd0def96c8b9c1c6efbf1f95130482a}.
\end{center}
Hence $\Omega_2(r)>\Omega_2(1350)$ for every $r\ge3001$.  The prefix and
half-line certificate overlap at $r=3001$.
\end{proof}

The target rational itself is regenerated from the wall recurrence.  It is
not embedded as a several-thousand-digit constant.  A one-percent upward
mutation of the target makes the constant coefficient of $H$ negative; this
checks that the tail gate can detect a false strengthening.

The half-line certificate is stronger than Proposition~\ref{prop:t2} needs,
and the surplus localizes the whole question.

\begin{corollary}[localization]\label{cor:localize}
The set of cells with $\Omega_t(r)<4/3$ is exactly
\[
 \bigl\{\,(r,2)\;:\;808\le r\le2339\,\bigr\},
\]
a contiguous window of $1532$ cells on the single stripe $t=2$.
\end{corollary}

\begin{proof}
The coefficient capacity of $H$ --- the largest target for which every
coefficient inequality still holds --- is
\[
 1.337695687131883021756995\ldots\;>\;\frac43 ,
\]
so the same construction at target $4/3$ retains eleven positive
coefficients and gives $\Omega_2(r)>4/3$ for every $r\ge3001$.  On the prefix
$2\le r\le3001$ the exact comparisons that prove Proposition~\ref{prop:t2}
also decide every cell against $4/3$: exactly $1532$ lie below it, they are
contiguous, and both boundaries are strict,
\[
 \Omega_2(807)>\tfrac43>\Omega_2(808),
 \qquad
 \Omega_2(2339)<\tfrac43<\Omega_2(2340).
\]
Every stripe $t\ge3$ is strictly above $4/3$: Proposition~\ref{prop:t3t4}
for $t=3,4$, and Theorem~\ref{thm:bulk} with Lemma~\ref{lem:transfer} for
$t\ge5$.
\end{proof}

The tail certificate is not tight at its own left edge, since
$\Omega_2(3001)=1.3394612629\ldots$; the capacity therefore lies below the
value it certifies against, as a valid bound must, and above $4/3$, as
Corollary~\ref{cor:localize} needs.  The extremal cell $(1350,2)$ lies inside
the window, so the entire sub-$4/3$ part of the array is one finite segment
of one stripe.

\section{The first higher stripes}\label{sec:firststripes}

The separator $4/3$ first appears on $t=3$.

\begin{proposition}\label{prop:t3t4}
For every integer $r\ge3$, $\Omega_3(r)>4/3$.  For every integer $r\ge4$,
$\rho(r,4)>4/3$ and consequently $\Omega_4(r)>4/3$.
\end{proposition}

\begin{proof}[Exact certificates]
For $t=3$, the bridge ratio is checked directly on $3\le r\le3001$ and a
degree-fourteen positive-coefficient polynomial proves the remaining
half-line.  Its coefficient capacity is
$1.4458344930865687354\ldots>4/3$; changing the target to $3/2$ creates
negative coefficients.

For $t=4$, the definition-level prefix checks the $997$ cells
$4\le r\le1000$.  At anchor $1000$, a degree-thirty-eight polynomial has all
$39$ coefficients strictly positive at target $4/3$.  Its coefficient
capacity is $1.3417349730931779905\ldots$, so the margin is thin but exact.
Lemma~\ref{lem:transfer} then gives the bridge statement.
\end{proof}

The two certificate joins are summarized in Table~\ref{tab:first}.
\begin{table}[ht]
\centering
\caption{Exact certificates for the first higher stripes.  Capacities are
reporting decimals of exact rationals.}\label{tab:first}
\begin{tabular}{@{}ccccc@{}}
\toprule
stripe & object & exact prefix & tail degree & capacity\\
\midrule
$t=3$ & $\Omega_3$ & $3\le r\le3001$ & 14 & $1.4458344930\ldots$\\
$t=4$ & $\rho(\cdot,4)$ & $4\le r\le1000$ & 38 & $1.3417349730\ldots$\\
\bottomrule
\end{tabular}
\end{table}

\section{The uniform \texorpdfstring{$4/3$}{4/3} barrier}\label{sec:bulk}

\begin{theorem}[bulk barrier]\label{thm:bulk}
For every pair of integers $j\ge t\ge5$,
\begin{equation}\label{eq:bulk}
 \rho(j,t)>\frac43.
\end{equation}
\end{theorem}

\begin{remark}
Theorem~\ref{thm:bulk} is used below only as a separator, but it is a
statement about the square-tail quotient in its own right.  The sharp global
constant of~\cite{squaretail} is $\min\rho=\rho(129,2)=1.0145761645\ldots$
over $2\le t\le r$; on the region $t\ge5$ the present theorem raises that to
$\rho>4/3$, a factor of $1.3141776\ldots$.  It does not claim that $4/3$ is
approached on any stripe; see \S\ref{sec:open}.
\end{remark}

The proof partitions the array by mechanisms rather than by a large bounding
box.  Two terminal identities and one finite core handle $j<504$.  Columns
$5$ through $10$ have exact positive-coefficient tails.  Above them, the
increasing wall level $\level_j$ divides the array into top and joint regions.

\subsection{Terminal cells and finite core}

The upper two coefficient cells are symbolic identities from
\cite{squaretail}:
\begin{equation}\label{eq:terminal}
 \rho(j,j)=\frac{j^2}{2},
 \qquad
 \rho(j,j-1)>\frac{2j^2}{9}.
\end{equation}
For a fixed column $t\ge5$, these give
\[
 \rho(t,t)\ge\frac{25}{2},
 \qquad
 \rho(t+1,t)>8.
\]

The remaining cells below row $504$ form the exact core
\begin{equation}\label{eq:corebox}
 7\le j\le503,\qquad 5\le t\le j-2.
\end{equation}
There are
$\sum_{j=7}^{503}(j-6)=123{,}753$ such cells.  A factorial-cancelled integer
formula tests
\[
 3N_{j,t}^2-4\bigl(2j^2\widehat S_{j,t}\widehat V_{j,t}\bigr)>0
\]
at every cell, which is exactly $\rho(j,t)>4/3$.  There are no failed or
degenerate cells.  The minimum inside the box occurs at $(j,t)=(503,5)$ and is
\[
 3.577109931649850292311076\ldots.
\]
The ordered comparison stream has digest
\begin{center}
\texttt{83797ef82eca71c6409a61f942bd81894d7352cf557c2bae2055bd40b9216402}.
\end{center}

\subsection{The fixed columns \texorpdfstring{$5\le t\le8$}{5 <= t <= 8}}

The unbounded certificate is uniform in form.  At an anchor $m_0$, exact
outward envelopes in a variable $x\ge0$ produce polynomials
$u_{\rm lo},p_{\rm hi},q_{\rm hi},v$.  The comparison polynomial is
\begin{equation}\label{eq:tailpoly}
 H_{t,m_0}(x)=u_{\rm lo}(x)^2
 -\frac43\,8p_{\rm hi}(x)q_{\rm hi}(x)v(x)^2.
\end{equation}
Strict positivity of every coefficient proves $\rho(j,t)>4/3$ for all
$j\ge m_0+1$.  A definition-level integer prefix reaches the anchor and also
checks the overlap cell $m_0+1$.  Table~\ref{tab:columns} gives the exact
joins.  The capacity is the largest target allowed simultaneously by all
coefficient inequalities, reported decimally only after the exact signs have
been decided.

\begin{table}[ht]
\centering
\caption{Complete fixed-column certificates at target $4/3$.}
\label{tab:columns}
\begin{tabular}{@{}ccccc@{}}
\toprule
$t$ & exact prefix & tail starts & degree & capacity\\
\midrule
5 & $5\le j\le1000$ & 1001 & 46 & $1.410970990331\ldots$\\
6 & $6\le j\le4006$ & 4007 & 54 & $1.518325000294\ldots$\\
7 & $7\le j\le12398$ & 12399 & 62 & $1.581611335023\ldots$\\
8 & $8\le j\le19999$ & 20000 & 70 & $1.420818931366\ldots$\\
\bottomrule
\end{tabular}
\end{table}

Several smaller anchors fail, which is useful evidence about the mechanism:
$m_0=2003$ fails for $t=6$ with capacity $1.00047\ldots$, and $m_0=6199$
fails for $t=7$ with capacity $1.16994\ldots$.  These failures do not weaken
the theorem; they locate where the chosen envelope becomes strong enough.

\subsection{The top region}

The top estimate already proved in~\cite{squaretail} has more room than its
original global anchor shows.  On
\[
 j\ge4,\qquad 2\le t\le j-2,\qquad \level_j\le(t-2)^2,
\]
its proof gives
\begin{equation}\label{eq:topchain}
 \rho(j,t)>
 \frac12\left(1-\frac2{4j-3}\right)^2
             \left(2-\frac1j\right)^2.
\end{equation}
Both factors increase with $j$.  Restricting to $j\ge504$ therefore yields
the exact floor
\begin{equation}\label{eq:topfloor}
 \rho(j,t)>
 \frac12\left(\frac{2011}{2013}\frac{1007}{504}\right)^2
 =\frac{4100936855929}{2058631521408}
 =1.9920693981\ldots>\frac43.
\end{equation}
No new asymptotic estimate is needed: the earlier, much smaller floor was the
same increasing product evaluated at $j=4$.

\subsection{The gate-compatible columns \texorpdfstring{$t=9,10$}{t=9,10}}

Columns $9$ and $10$ use the top region below an anchor and a
positive-coefficient tail above it.  Since $\level_j$ is increasing, the
single inequality $\level_{m_0}\le(t-2)^2$ places every row
$504\le j\le m_0$ in \eqref{eq:topfloor}.  Table~\ref{tab:ceiling} records the
two joins.

\begin{table}[ht]
\centering
\caption{Prefix-free joins for columns $9$ and $10$.  The finite core covers
$j<504$, the top theorem covers through $m_0$, and the tail begins at
$m_0+1$.}\label{tab:ceiling}
\begin{tabular}{@{}cccccc@{}}
\toprule
$t$ & $m_0$ & certified gate & tail starts & degree & capacity\\
\midrule
9 & 68133 & $\level_{68133}<49$ & 68134 & 78 & $1.660681546638\ldots$\\
10 & 200000 & $\level_{200000}<64$ & 200001 & 86 & $1.515648188324\ldots$\\
\bottomrule
\end{tabular}
\end{table}

For $t=9$ the splice is exact at the level crossing: outward intervals prove
\[
 \level_{68133}<49<\level_{68134}.
\]
Thus $68133$ is the largest top-compatible anchor.  A smaller build at
$55000$ fails, while the ceiling build passes with every one of its $79$
coefficients positive.  This corrects a tempting but false diagnostic: the
envelope route is not dead for $t\ge9$; it is inadequate only at anchors that
sit too low.  For $t=10$, the build at $100000$ fails and the build at
$200000$ passes.  A separate sharpened-corridor Sturm certificate also closes
the $t=10$ joint range, but it is not used in the assembly.
For $t=9$, the sharpened-corridor slices combined with
Proposition~\ref{prop:dlaw} give a second certificate chain across the dip
$49\le\level_j\le768$.  This chain is likewise corroborating rather than
load-bearing.

\subsection{The joint region for \texorpdfstring{$t\ge11$}{t >= 11}}

It remains to prove the $4/3$ floor when
\begin{equation}\label{eq:jointdomain}
 j\ge504,\qquad t\ge11,\qquad \level_j>(t-2)^2.
\end{equation}
The corridor reduction in~\cite{squaretail} uses a real parameter $Y$ and the
rational function
\begin{align}
 G&=2t(2t+1),&
 \alpha&=\frac{(t-1)(2t-1)}{t(2t+1)},&
 \beta&=\frac{t(2t+1)}{(t+1)(2t+3)},\nonumber\\
 x_0&=\frac{Y-5}{G},&
 u_0&=\alpha\frac{Y-5}{Y},&
 R_0&=1+x_0+(x_0u_0+x_0^2)\frac{1-\beta}{1-u_0}.
\label{eq:R0}
\end{align}

\begin{lemma}[joint core]\label{lem:jointcore}
For every integer $t\ge11$ and every real
$Y\ge(t-2)(t-7)$, putting $a=Y+5(t-2)$ gives
\begin{equation}\label{eq:core2720}
 \frac{[(2t-1)R_0-1]^2}{2aR_0}\ge\frac{27}{20}.
\end{equation}
\end{lemma}

\begin{proof}[Exact polynomial certificate]
Clear the positive denominators in \eqref{eq:core2720}.  For
$t=11,\dots,16$, the resulting univariate polynomial after
$Y=(t-2)(t-7)+Q$ is positive on $Q\ge0$ by exact Sturm sequences: the sign
variations at $0$ and $+\infty$ agree and the value at $0$ is positive.

For every $t\ge17$, substitute
\[
 t=S+17,\qquad Y=Q+(S+15)(S+10).
\]
The cleared bivariate polynomial has $105$ monomials in $(Q,S)$.  Every
coefficient is nonnegative and the constant coefficient is positive.  It is
therefore positive on the whole quadrant $Q,S\ge0$.  The indexed coefficient
digest is
\begin{center}
\texttt{9726005502a811c97584bc22b98eeced6577d4b18a5c3c45df3eb178239c69b4}.
\end{center}
The two certificates cover all integers $t\ge11$ and the full real half-line
in $Y$.
\end{proof}

The corridor chain also gives, on \eqref{eq:jointdomain},
\begin{equation}\label{eq:attached}
 \rho(j,t)\ge\frac{27}{20}
 \left(1-\frac2{4j-3}\right)^2
 \left(1-\frac2j\right)^2.
\end{equation}
The two attached factors increase with $j$.  At $j=504$ their product is
\[
 \frac{2011}{2013}\frac{251}{252}
 =\frac{504761}{507276}.
\]
Consequently
\begin{equation}\label{eq:jointfloor}
 \rho(j,t)\ge
 \frac{27}{20}\left(\frac{504761}{507276}\right)^2
 =\frac{254783667121}{190614029760}
 =1.3366469794\ldots>\frac43.
\end{equation}
The row-free core needed to clear $4/3$ after these attached factors is exactly
\[
 \frac{4/3}{(504761/507276)^2}
 =\frac{343105253568}{254783667121}
 =1.3466532507\ldots,
\]
which is strictly below $27/20$.

To verify that \eqref{eq:jointdomain} lies in the domain of
Lemma~\ref{lem:jointcore}, the corridor gives
$Y\ge\level_j-5(t-2)$.  Hence
\[
 Y>(t-2)^2-5(t-2)=(t-2)(t-7).
\]
Moreover
\begin{equation}\label{eq:dichotomyidentity}
 (t-2)^2-5t=(t-9)^2+9(t-9)+4>0
\end{equation}
for every integer $t\ge9$, so the original joint hypothesis
$\level_j>5t$ also holds.

\subsection{Coverage proof}

\begin{proof}[Proof of Theorem~\ref{thm:bulk}]
Fix $j\ge t\ge5$.  If $j=t$ or $j=t+1$, use
\eqref{eq:terminal}.  If $t+2\le j\le503$, use the finite core
\eqref{eq:corebox}.  Thus assume $j\ge504$.

For $5\le t\le8$, Table~\ref{tab:columns} closes the entire column.  For
$t=9,10$, Table~\ref{tab:ceiling} and the top floor
\eqref{eq:topfloor} close the column below the tail, while
\eqref{eq:tailpoly} closes it above the anchor.

Finally let $t\ge11$.  If $\level_j\le(t-2)^2$, use the top floor.  Otherwise
\eqref{eq:jointdomain} holds, and the joint floor \eqref{eq:jointfloor}
applies.  These two inequalities are complementary, with equality assigned
to the top side, so no cell is omitted.  This exhausts the domain.
\end{proof}

\section{Assembly and sharpness}\label{sec:assembly}

\begin{proof}[Proof of Theorem~\ref{thm:main}]
Proposition~\ref{prop:t2} proves
$\Omega_2(r)\ge\Omw$, with equality only at $r=1350$.
Proposition~\ref{prop:t3t4} proves $\Omega_t(r)>4/3$ for $t=3,4$.
For $t\ge5$, Theorem~\ref{thm:bulk} and Lemma~\ref{lem:transfer} give
\[
 \Omega_t(r)\ge\min_{t\le j\le r}\rho(j,t)>\frac43.
\]
The exact comparison \eqref{eq:separator} now gives \eqref{eq:main} and shows
that no equality is possible away from $(1350,2)$.  Multiplying by $2$ proves
the raw formulation \eqref{eq:rawsharp}; equality at the stated cell shows
that its constant cannot be increased.
\end{proof}

The master join rebuilds $\Omw$, checks its target fingerprint and its exact
comparison with $4/3$, validates the three closed lower-stripe records, every
regional certificate and splice, and the symbolic coverage arithmetic.  It
has $22$ named gates; all $22$ pass and there are no conditional legs.

\subsection{Why monotonicity is not the proof}

It is natural to guess that $\Omega_{t+1}(r)\ge\Omega_t(r)$ pointwise.  That
statement is false.  An exact integer comparison at row $45082$ gives
\begin{equation}\label{eq:monotonicityfail}
 \Omega_3(45082)<\Omega_2(45082).
\end{equation}
The witness is the exact integer
\[
 D=L(G)_{r-3}L(C)_{r-2}-L(G)_{r-2}L(C)_{r-3},
 \qquad
 G=(1+z)M_{2r-2},\quad C=M_{2r-1},
\]
whose sign is that of $\Omega_3(r)-\Omega_2(r)$ because $L(C)_q>0$ on this
range.  At $r=45082$ it is negative.  Its magnitude has bit length
$5{,}416{,}154$ and $1{,}630{,}425$ decimal digits, and the SHA-256 digest of
its big-endian byte encoding is
\begin{center}
\texttt{0aed686f60e5eaf6056db6cb8e1f59aa99610dc26293b7f9f57fb53d2f8ed47e}.
\end{center}
Only the top few coefficients of each wall row are consumed, so the companion
producer carries a fixed-width top window of the recurrence and reaches
$r=45082$ in well under a minute.  It checks that window against the full
definition on $5\le r\le40$ and rejects a recurrence mutation before
rebuilding the displayed witness.
This does not threaten Theorem~\ref{thm:main}; it only rules out a shortcut.
The transfer lemma and the $4/3$ barrier compare each stripe with a common
rational floor instead.

The failed low-anchor tail builds have the same status.  They identify proof
loss in a particular envelope, not counterexamples to the theorem.  In
particular, the passing build at the largest top-compatible anchor for $t=9$
shows why a global statement that the envelope route ``dies at $t=9$'' would
be incorrect.

\section{A decrement shell and the open profile}\label{sec:open}

The bulk proof exposes a quantitative law that is stronger than the corridor
bound needed for Theorem~\ref{thm:main}.  In the deposited ladder indexing,
put
\[
 X_{m,k}=2k(2k+1)\frac{B_k(m)}{B_{k-1}(m)},
 \qquad d_{m,k}=X_{m,k-1}-X_{m,k}.
\]

\begin{proposition}[certified decrement shell]\label{prop:dlaw}
For every integer $m\ge68133$ with $\level_{m+1}\le768$, and every
$k\in\{5,\ldots,11\}$,
\begin{equation}\label{eq:dlaw}
 (k-1)d_{m,k}\le\frac65\sqrt{\level_{m+1}}.
\end{equation}
\end{proposition}

\begin{proof}[Exact two-leg certificate]
The first leg propagates the normalized $B$--$C$ ladder with
$192$-bit outward fixed-point intervals.  It checks
\eqref{eq:dlaw} at every row $60000\le m\le262144$ and all seven indices,
with no failure, and deposits the complete interval state at the last row.

For the second leg, write $\beta_k=(\pi/2)B_k$ and use the polynomial flow
\[
 \beta_k'=\gamma_k,\qquad \gamma_k'=\beta_{k-1}.
\]
The connection polynomials are the exact closure of this flow, so one
interval anchor suffices.  Injection errors are majorized in the shifted
variable $x=\tau-\tau_{262144}$; the logarithmic tail moments are bounded
exactly by
\[
 \sum_{j>m}\frac{\bigl(\tfrac12\log(j/m)\bigr)^i}{j^2}
 \le \frac{i!}{2^i m}+\frac{2(i/8)^i}{m^2}.
\]
Fifty-seven contiguous windows of width $7/20$ cover
$0\le x\le399/20$.  On each window, rational interval subdivision proves
the seven cleared inequalities and certifies that every lower envelope
which is squared is nonnegative.  At the right endpoint the lower bound for
$\level$ exceeds $768$; strict increase of $\level_m$ then closes the stated
window.  The minimum reported relative margin rises from
$0.2719\ldots$ to $0.6261\ldots$; these decimals report exact rational
signs.  Retargeting the same certificate from $6/5$ to $4/5$ fails in the
first window at six indices, providing a negative control.
\end{proof}

Proposition~\ref{prop:dlaw} is not needed for the sharp constant: the
top--tail join already closes $t=9$.  Combined with the sharpened-corridor
slices, however, it gives a second separately gated proof chain for that
stripe, parallel to the two chains available for $t=10$.  These routes share
the wall definitions and upstream results from~\cite{squaretail}; they are
not claimed to be fully independent implementations.

The open object is now the profile behind \eqref{eq:dlaw}, not the displayed
inequality itself.  Exact data suggest a scaling law of the form
\[
 d_{m,k}\sim\sqrt{\level_{m+1}}\,
 g\!\left(\frac{k-1}{\sqrt{\level_{m+1}}}\right),
\]
with the measured normalized decrement near $0.88$ on the present window and
drifting toward $1$.  Determining $g$, proving the limiting constant, and
turning the profile into a uniform-in-$t$ deficit transfer remain open.
So do the individual higher-stripe minima: the proof uses $4/3$ as a
separator, not as their sharp value.  Finally, the present results are exact
computer-assisted theorems rather than kernel-checked theorems; a
proof-assistant replay remains a separate program.

\section{Reproducibility and claim boundary}\label{sec:repro}

The proof uses symbolic arguments and exact computation in distinct roles.
Table~\ref{tab:claims} states the boundary explicitly.

\begin{table}[ht]
\centering
\caption{Status of the durable claims.}\label{tab:claims}
\small
\begin{tabularx}{\textwidth}{@{}P{0.25\textwidth}P{0.22\textwidth}Y@{}}
\toprule
claim & status & anchor\\
\midrule
Theorem~\ref{thm:main} & exact computer-assisted theorem & written transfer
and coverage proof plus the $22$-gate exact join\\
terminal and top formulas & all-parameter theorems & symbolic derivations in
the text and in~\cite{squaretail}\\
finite core and prefixes & finite exact theorems on stated ranges & integer
comparison streams with counts and digests\\
tail and joint half-lines & exact computer-assisted theorems & rational
coefficient signs and exact Sturm sequences\\
reported decimal capacities & derived reporting values & exact rationals in
the certificate records\\
individual higher-stripe minima & open & exploratory numerical census only\\
decrement law \eqref{eq:dlaw} & exact computer-assisted theorem & $192$-bit
finite ladder, $57$ exact envelope windows, and rejected $4/5$ control\\
limiting decrement profile & open & exploratory scaling data only\\
\bottomrule
\end{tabularx}
\end{table}

No floating-point comparison decides a theorem gate.  The finite core uses
exact integers; prefix minima use exact cross multiplication; half-line tails
use rational interval envelopes and coefficientwise positivity; the joint
certificate uses integer polynomials and exact Sturm signs.  Every reported
decimal is obtained after the exact verdict.

One inherited interval constructor deserves disclosure.  Its nominal upper
endpoint for $\pi$ used truncation and therefore produced a singleton lying one
fixed-point unit too low.  The replay first detects the defect and then repairs
the interval by an outward ceiling.  Every current consumer immediately scales
the interval by $1/2$; exact integer rounding maps the defective and corrected
cells to the same interval for $\pi/2$.  Thus the repair restores the source
contract while changing every downstream coefficient by exactly zero.  A
hard-coded $100$-decimal-digit enclosure for $\pi$ rejects the old constructor
and accepts the repaired one.

The release record contains the exact certificate records, the ordered
comparison digest and summary extrema for every prefix stream, the
positive-polynomial certificates, the failed-anchor controls, the exact
crossing records, and the master join, together with the producer that
regenerates each stream in full.  The reader-facing source map names the
load-bearing artifacts and their digests.  This paper makes no kernel-verified
claim.  The manuscript and its companion replay carry the identifier
\href{https://doi.org/10.5281/zenodo.21866366}{doi:10.5281/zenodo.21866366}.

\section{Conclusion}

The bridge constant is controlled by a single exceptional cell, but proving
that fact requires a global argument.  The direct $t=2$ certificate finds the
cell.  The lossless transfer converts the rest of the question into a common
$4/3$ square-tail barrier.  Terminal identities, a finite core, fixed columns,
and the top--joint dichotomy then cover every remaining row without an
asymptotic gap.

This separates two constants that play different roles.  The factor $2$ in
the earlier bridge theorem is structural and easy to transport through the
Tur\'an triangle.  The number
$2.655018313913361199726406175676\ldots$ is extremal: it is the largest raw
factor valid over the whole shift wall, and its equality cell is unique.

\section*{Acknowledgments}

This work was carried out independently, without institutional affiliation or
funding.

Large language models were used as mathematical interlocutors: to test proof
routes, carry out derivations, and organize exact certificate searches.  Every
mathematical claim in the paper is supported by a written proof or an exact
certificate, and responsibility for the result is mine.

\begin{thebibliography}{9}

\bibitem{delannoy}
J.~Councilman,
\emph{Strict total positivity of a real Delannoy multiplication table: the
sharp half-line threshold}, v2.0, Zenodo, 2026;
\href{https://doi.org/10.5281/zenodo.21778524}{doi:10.5281/zenodo.21778524}.

\bibitem{greenpath}
J.~Councilman,
\emph{green-path-dpp: exact rational determinantal algorithms for grounded
paths}, v0.1.0, Zenodo, 2026;
\href{https://doi.org/10.5281/zenodo.21778518}{doi:10.5281/zenodo.21778518}.

\bibitem{squaretail}
J.~Councilman,
\emph{Odd cycles, sech spectra, and a square-tail inequality}, v0.1.0,
Zenodo, 2026;
\href{https://doi.org/10.5281/zenodo.21853982}{doi:10.5281/zenodo.21853982}.

\end{thebibliography}

\end{document}
